\documentclass[11pt]{article}
\usepackage[T1]{fontenc}
\usepackage[margin=1in]{geometry}
\usepackage{enumitem}
\usepackage{hyperref}
\setlength{\parindent}{1.5em}
\setlength{\parskip}{6pt}
% =====================================================================
%  EDIT YOUR DETAILS HERE — this ONE file feeds every abstract & deck.
%  (Change a value, recompile; all 36 documents update.)
% =====================================================================
\newcommand{\seminarName}{Aflah Muhammed P}      % <-- your full name (guessed from your email; please verify)
\newcommand{\seminarRoll}{TVE00CS000}            % <-- your university register/roll number
\newcommand{\seminarClassRoll}{00}               % <-- your class roll no (the short one)
\newcommand{\seminarGuide}{Prof.\ <Guide Name>}  % <-- your guide
\newcommand{\seminarDept}{Department of Computer Science and Engineering}
\newcommand{\seminarCollege}{College of Engineering Trivandrum}
\newcommand{\seminarShortCollege}{CET}           % shown in the slide footer, e.g. "Aflah Muhammed P (CET)"
\newcommand{\seminarShortTitle}{B.Tech Seminar}  % middle of the slide footer
\newcommand{\seminarDate}{July 31, 2026}
% ---------------------------------------------------------------------
% Optional: put a logo image named  cet_logo.png  in this folder and it
% will appear on every title slide automatically. If absent, it is skipped.
% =====================================================================

\begin{document}
\begin{center}
{\LARGE\bfseries AgentRewardBench: Evaluating Automatic Evaluators of Web Agents}\\[1.1em]
{\large\bfseries Seminar Abstract}\\[0.6em]
{\normalsize \seminarDate}
\end{center}
\vspace{0.6em}
\section*{Abstract}
Autonomous web agents built on large language models can navigate real websites to complete multi-step tasks such as shopping, form filling, and information retrieval. Measuring whether an agent actually succeeded, however, is deceptively difficult. Benchmarks typically rely on rule-based checks that compare final page states or exact strings, which are brittle and reject valid alternative solutions. Manual expert review is accurate but does not scale to the thousands of trajectories modern research produces. Increasingly, practitioners deploy large language models as automatic judges, yet the reliability of these evaluators has itself gone largely unverified, leaving a critical blind spot in how progress on web agents is measured and rewarded.

This paper introduces \emph{AgentRewardBench}, among the first benchmarks designed to evaluate the evaluators of web-agent trajectories. It comprises 1,302 trajectories spanning five environments (WebArena, VisualWebArena, WorkArena L1 and L2, and AssistantBench) generated by four agent language models, each reviewed by an expert who annotates success, side effects, and repetitive behaviour. Using these gold labels, the authors benchmark twelve LLM judges alongside conventional rule-based checks, reporting precision and recall against expert judgments. Two findings stand out: no single LLM judge performs best across all benchmarks, and the widely used rule-based evaluation systematically under-reports agent success. The seminar walks through the dataset construction, the taxonomy of evaluators, the quantitative results, and their implications for building more flexible and trustworthy automatic evaluation of web agents.
\section*{Key Reference Papers}
\begin{enumerate}
  \item X. H. L\`u, A. Kazemnejad, N. Meade, A. Patel, D. Shin, A. Zambrano, K. Sta\'nczak, P. Shaw, C. J. Pal, and S. Reddy, ``AgentRewardBench: Evaluating Automatic Evaluations of Web Agent Trajectories,'' arXiv:2504.08942, 2025.
  \item S. Zhou, F. F. Xu, H. Zhu, et al., ``WebArena: A Realistic Web Environment for Building Autonomous Agents,'' in Proc. ICLR, 2024, arXiv:2307.13854.
  \item J. Y. Koh, R. Lo, L. Jang, et al., ``VisualWebArena: Evaluating Multimodal Agents on Realistic Visual Web Tasks,'' arXiv:2401.13649, 2024.
\end{enumerate}
\vspace{0.8em}
\noindent\textbf{Submitted By}\\
\seminarName\\
\seminarRoll

\vspace{0.6em}
\noindent\textbf{Guide}\\
\seminarGuide
\end{document}
